\documentclass[a4paper,12pt,headsepline]{report}

%----------------- PDF CONFIG ----------------- %
\pdfinfo{    
     /Title (PDF-Titel) 
     /Subject   (PDF-Thema)    
     /Author  (Vorname Nachname) 
     /Keywords   (Stichwort1,Stichwort2)      
} 

\title{MyTitle}
\author{theAuthor}
\date{1.1.2000}



%----------------- PAKETE INKLUDIEREN ----------------- %

\usepackage{geometry} % Packet für Seitenrandabständex und Einstellung für Seitenränder
\usepackage[ngerman]{babel} % deutsche Silbentrennung

\usepackage{booktabs} %entzerrt die Tabellenzeilen und bietet verschieden dicke Unterteilungslinien
\usepackage{longtable} % Tabellen können sich nicht über mehrere Seiten 
\usepackage{graphicx} % kann LaTeX Grafiken einbinden

\usepackage[applemac]{inputenc} % Umlaute unter Mac werden automatisch gesetzt
\usepackage[T1]{fontenc} % Zeichenencoding
\usepackage{lmodern} % typographische Qualität 
\frenchspacing % Schaltet den zusätzlichen Zwischenraum ab
\usepackage{fix-cm}
\usepackage{hyperref} % verwandelt alle Kapitelüberschriften, Verweise aufs Literaturverzeichnis und andere Querverweise in PDF-Hyperlinks
\usepackage{color}
\usepackage{url}


\usepackage[nottoc]{tocbibind}



% für Listings
\usepackage{listings}
\lstset{numbers=left, numberstyle=\tiny, numbersep=5pt, stepnumber=4, keywordstyle=\color{black}\bfseries\itshape, stringstyle=\ttfamily,showstringspaces=false,basicstyle=\footnotesize,captionpos=b}
\lstset{language=java}



%----------------- FARBEN DEFINIEREN ----------------- %
\definecolor{gray}{gray}{0.95} % Listingsbackground

%----------------- LAYOUT SETZEN ----------------- %
\geometry{left=2cm, right=2cm, top=2.5cm, bottom=2cm}
\linespread {1.25}\selectfont %1.25 da er von Haus aus 1.2 ist und 1,25 * 1,2 = 1,5 isch




%-------##-------##-------##------- ANFANG INHALT -------##-------##-------##-------%
\begin{document}



%\pagenumbering{roman} % Seitennummer

%----------------- DECKBLATT -----------------%
 %----------------- KONFIGURATION ----------------- %
\pagestyle{empty} % enthalten keinerlei Kopf oder Fuß 

%----------------- HDA FBI Logo ----------------- %
\begin{figure}[t]
	\centering
	\includegraphics[width=0.6\textwidth]{Abb/logo_fbi}
\end{figure}

%----------------- INHALT ----------------- %

\begin{center}
\Large Hochschule Darmstadt \\
\normalsize \textsc{- Fachbereich Informatik -} \\

% Whitespace
\vspace{105 pt}

\Huge Name der Arbeit \\ 
\normalsize
\vspace{20 pt}

Abschlussarbeit zur Erlangung des akademischen Grades \\ 
Bachelor of Science (B.Sc.) 

\vspace{75 pt}


vorgelegt von \\
\vspace{5 pt}
Can Kedik \\
731620
\vspace{115 pt}

\begin{tabular}[h]{p{4cm}l l}
	Referent: & Prof. Dr. Ronald C. Moore\\
	Korreferent: & Prof. Dr. Bettina Harriehausen-M"uhlbauer
\end{tabular}


\end{center}

 
%----------------- ABSTRACT -----------------%
 %----------------- KONFIGURATION ----------------- %
\pagestyle{empty} % enthalten keinerlei Kopf oder Fuß


\chapter*{Abstract} % (fold)
\label{cha:abtract}
	Lorem ipsum dolor sit amet, consetetur sadipscing elitr, sed diam nonumy eirmod tempor invidunt ut labore et dolore magna aliquyam erat, sed diam voluptua. At vero eos et accusam et justo duo dolores et ea rebum. Stet clita kasd gubergren, no sea takimata sanctus est Lorem ipsum dolor sit amet. Lorem ipsum dolor sit amet, consetetur sadipscing elitr, sed diam nonumy eirmod tempor invidunt ut labore et dolore magna aliquyam erat, sed diam voluptua. At vero eos et accusam et justo duo dolores et ea rebum. Stet clita kasd gubergren, no sea takimata sanctus est Lorem ipsum dolor sit amet.

\begin{quote}
Dies ist ein Zitat.
\end{quote}

verstand, scheinen nun doch vorueber zu Dies ist der Text sein. \\
siehe: http://janeden.net/die-praeambel

 
 
%----------------- VERZEICHNISSE -----------------%
\tableofcontents % Inhaltverzeichnis

% ----- Abbildungen ----- %
%\addcontentsline{toc}{section}{Abbildungsverzeichnis} % falls in Inhalsverzeichnis
\listoffigures

% ----- Tabellen----- %
% \addcontentsline{toc}{section}{Tabellenverzeichnis}  % falls in Inhalsverzeichnis
% \fancyhead[L]{Abbildungsverzeichnis / Abkürzungsverzeichnis} %Kopfzeile links
\listoftables

% ----- Listings ----- %
% Listingverzeichnis soll im Inhaltsverzeichnis auftauchen
% \addcontentsline{toc}{section}{Listingverzeichnis}
% \fancyhead[L]{Abbildungs- / Tabellen- / Listingverzeichnis} %Kopfzeile links
\renewcommand{\lstlistlistingname}{Listingverzeichnis}
\lstlistoflistings

\pagestyle{plain} % zurueck setzen von roemische seitenanzahl


%----------------- KAPITEL : EINLEITUNG  ----------------- %
\pagenumbering{arabic}
\chapter{Einf"uhrung}
\label{cha:einfuehrung}

Lorem ipsum dolor sit amet, consetetur sadipscing elitr, sed diam nonumy eirmod tempor invidunt ut labore et dolore magna aliquyam erat, sed diam voluptua. At vero eos et accusam et justo duo dolores et ea rebum. Stet clita kasd gubergren, no sea takimata sanctus est Lorem ipsum dolor sit amet. Lorem ipsum dolor sit amet, consetetur sadipscing elitr, sed diam nonumy eirmod tempor invidunt ut labore et dolore magna aliquyam erat, sed diam voluptua. At vero eos et accusam et justo duo dolores et ea rebum. Stet clita kasd gubergren, no sea takimata sanctus est Lorem ipsum dolor sit amet.



%----------------- KAPITEL : IDEE  ----------------- %	
\chapter{Idee} %fällt wahrscheinlich weg

%----------------- KAPITEL : FRAMEWORK  ----------------- %
\chapter{Das Framework}
\section{Verwendete Tools}
\section{Aufbau der Software}
\section{Einsatz im Netzwerk} 

%----------------- KAPITEL : TESTAUFBAU  ----------------- %
\chapter{Testaufbau}
\section{Pis Real}
\section{Pis Virtuell}

%----------------- KAPITEL : VERSUCHE  ----------------- %
\chapter{Versuche}
\section{Raspberry Pi Versuche}
\subsection{Gemachte Tests}
\subsection{Ergebnisse}
\section{Virtual Machine Versuche}
\subsection{Gemachte Tests}
\subsection{Ergebnisse}

%----------------- KAPITEL : VERGLEICH------------ %	
\chapter{Vergleich VM/HW}
\section{Versuchsergebnisse}
\section{Kosten nutzen Faktor}

%----------------- KAPITEL : FAZIT  ----------------- %	
\chapter{Fazit}

%----------------- KAPITEL : ANHANG  ----------------- %	


\end{document}
